\documentclass[12pt]{article}

\usepackage{sbc-template}
\usepackage{graphicx}
\usepackage{grffile}
\usepackage{url}
\usepackage{booktabs}
\usepackage{multirow}
\usepackage{tabularx}
\usepackage{caption}
\usepackage{float}

\title{Miniservidor HTTP/1.1 com Resolução DNS Local sobre TCP e R-UDP}

\author{Brenda Evilly Oliveira de Moura}

\address{%
  Universidade Federal do Piauí (UFPI)\\
  Campus Senador Helvídio Nunes de Barros (CSHNB)\\
  Bacharelado em Sistemas de Informação\\
  Picos -- PI -- Brasil\\
  \email{brenda.moura@ufpi.edu.br}
}

\begin{document}

\maketitle

\begin{abstract}
  Este relatório descreve a evolução de um sistema de transferência de arquivos
  para um miniservidor web capaz de atender requisições HTTP/1.1 sobre TCP nativo
  e sobre uma camada R-UDP desenvolvida previamente. O cliente web passou a
  depender de um módulo DNS simplificado em UDP antes de qualquer conexão ao
  servidor. Foram conduzidos experimentos em ambiente Docker com três cenários
  de rede simulados por \emph{tc netem}, comparando vazão, tempo de resposta e
  taxa de erro entre os modos de transporte. Os resultados indicam desempenho
  superior e maior estabilidade do TCP, enquanto o R-UDP apresentou limitações
  marcantes sob perda de pacotes, especialmente em arquivos maiores que 100~kB.
\end{abstract}

\section{Introdução}

A World Wide Web consolidou o HTTP como protocolo de aplicação dominante para
distribuição de documentos hipertexto e recursos estáticos \cite{rfc2616,rfc9110}.
Antes de qualquer requisição HTTP, o cliente precisa resolver um nome simbólico
para um endereço IP, função tradicionalmente desempenhada pelo DNS sobre UDP
\cite{rfc1035}. Na pilha TCP/IP, o UDP oferece entrega sem confirmação
\cite{rfc768}, enquanto o TCP garante confiabilidade, controle de fluxo e
ordenação \cite{rfc793,tanenbaum2011}.

Na disciplina de Redes de Computadores II, o trabalho anterior implementou
transferência de arquivos por meio de um protocolo textual customizado, com
alternância entre TCP e R-UDP (Stop-and-Wait) em contêineres Docker. A presente
avaliação estende essa base para simular navegação web: o cliente consulta um
servidor DNS local, obtém o IP do \emph{host} web e então realiza requisições
GET HTTP/1.1, servindo páginas HTML e arquivos binários de tamanhos distintos.

O objetivo deste documento é registrar a arquitetura implementada, a metodologia
de testes com emulação de atraso e perda, os resultados quantitativos obtidos
em 240 execuções planejadas e a validação do fluxo DNS $\rightarrow$ HTTP por
meio de capturas Wireshark. As análises buscam responder como a perda simulada
afeta o DNS e o download HTTP, qual o impacto dos cabeçalhos HTTP frente ao
protocolo textual anterior e se a sequência observada na rede corresponde à
ordem esperada da arquitetura em camadas \cite{kurose2017,carpenter2010}.

\section{Fundamentação Teórica}

O HTTP/1.1 define mensagens de requisição e resposta compostas por linha inicial,
cabeçalhos e corpo opcional \cite{rfc2616}. Campos como \texttt{Host},
\texttt{Content-Type} e \texttt{Content-Length} permitem identificar o recurso,
declarar o tipo de mídia e informar o tamanho do payload. No contexto acadêmico
deste projeto, incluiu-se ainda o cabeçalho personalizado \texttt{X-Custom-Auth},
calculado como hash SHA-256 da matrícula concatenada ao nome do estudante.

O DNS padrão utiliza consultas e respostas UDP na porta 53, com registros do
tipo A mapeando nomes para endereços IPv4 \cite{rfc1035}. Por operar sem
retransmissão nativa, consultas perdidas dependem de temporizadores e novas
tentativas na camada de aplicação \cite{kurose2017,tanenbaum2011}. O TCP, por
sua vez, realiza \emph{handshake}, transferência orientada a fluxo de bytes e
encerramento ordenado \cite{rfc793}.

A camada R-UDP reutilizada neste trabalho implementa confiabilidade Stop-and-Wait
sobre UDP, com cabeçalho binário contendo autenticação, \emph{checksum} e
números de sequência. Cada segmento de dados aguarda confirmação antes do envio
do próximo, o que simplifica a implementação, mas limita a vazão quando o RTT
aumenta ou quando há perda \cite{tanenbaum2011}. A emulação de canal foi feita
com \texttt{tc netem}, ferramenta consolidada para introduzir atraso e perda em
laboratório \cite{netem2005}.

\section{Metodologia}

Esta seção descreve a arquitetura experimental, os protocolos de aplicação
(DNS, HTTP e R-UDP), os cenários de rede simulados e o procedimento adotado
para coleta, análise e validação cruzada dos dados.

\subsection{Arquitetura experimental}

O ambiente de testes foi conteinerizado com Docker Compose \cite{docker2024},
utilizando três serviços baseados em Ubuntu 22.04 conectados à rede bridge
\texttt{172.29.0.0/24}. O contêiner DNS (\texttt{redes2-t02-dns}) ocupa o
endereço \texttt{172.29.0.5} e escuta consultas UDP na porta 53. O contêiner
servidor web (\texttt{redes2-t02-server}) responde em \texttt{172.29.0.10},
expondo HTTP/TCP na porta 8080 e HTTP/R-UDP na porta 8081. O contêiner cliente
(\texttt{redes2-t02-client}), em \texttt{172.29.0.20}, inicia com a
\emph{capability} \texttt{NET\_ADMIN} para aplicar regras \texttt{tc} em sua
própria interface \texttt{eth0}.

A aplicação foi desenvolvida em Python 3, sem bibliotecas de manipulação de
pacotes de alto nível. O cliente web \textbf{não} recebe o IP do servidor como
parâmetro obrigatório de conexão: ele resolve \texttt{www.redes.local} junto
ao DNS e só então abre a sessão HTTP no endereço retornado. Os arquivos servidos
incluem \texttt{index.html} (624~bytes), \texttt{arquivo\_100kb.bin},
\texttt{arquivo\_1mb.bin} e \texttt{arquivo\_10mb.bin}, gerados pelo script
\texttt{generate\_www\_files.py} antes de cada sessão de testes.

As métricas de cada transferência (tempo de DNS, tempo HTTP, duração total,
bytes recebidos, vazão, retransmissões R-UDP e indicador de sucesso) foram
registradas pela classe \texttt{WebTransferMetrics} no arquivo
\texttt{logs/transfers.jsonl}, no formato JSON Lines. A degradação de rede é
introduzida no \emph{egress} do cliente: todo tráfego DNS e HTTP atravessa o
\texttt{netem} configurado em \texttt{eth0} antes de alcançar os demais nós.

\subsection{Protocolo DNS simplificado}

O módulo DNS implementado neste trabalho não segue integralmente o RFC~1035, mas
reproduz consultas tipo A em formato compacto. A consulta UDP contém: ID
(2~bytes), tamanho do nome (2~bytes) e nome em UTF-8. A resposta repete ID e
nome e acrescenta quatro octetos IPv4; o endereço \texttt{0.0.0.0} indica falha
de resolução. A zona é lida do arquivo \texttt{dns/hosts.txt}, mapeando
\texttt{www.redes.local} para \texttt{172.29.0.10}.

O cliente DNS utiliza \texttt{DNS\_TIMEOUT} de 1,0~s e até três retransmissões
na camada de aplicação. O temporizador efetivo varia conforme o cenário: 0,5~s
(A), 1,0~s (B) e 2,0~s (C). O tempo de resolução é medido com
\texttt{time.perf\_counter()} entre o primeiro envio e a recepção da resposta
válida.

\subsection{Protocolo HTTP/1.1}

Após a resolução DNS, o cliente monta uma requisição textual HTTP/1.1, por
exemplo \texttt{GET /index.html HTTP/1.1}, seguida de \texttt{Host},
\texttt{X-Custom-Auth} e linha em branco. O servidor valida o hash de
autenticação, localiza o recurso em \texttt{www/} e responde com
\texttt{HTTP/1.1 200 OK} ou \texttt{404 Not Found}, incluindo
\texttt{Content-Type}, \texttt{Content-Length} e novamente
\texttt{X-Custom-Auth}.

No modo TCP, a requisição é enviada após \texttt{connect()} à porta 8080 do IP
resolvido; a resposta é lida até o encerramento da conexão. No modo R-UDP, a
mensagem HTTP completa (requisição e resposta) trafega como \emph{payload} dos
pacotes Stop-and-Wait na porta 8081, preservando os mesmos cabeçalhos de
aplicação.

\subsection{Protocolo R-UDP sobre HTTP}

O R-UDP reutiliza o cabeçalho binário de 52~bytes fixos do trabalho anterior
(\emph{magic}, versão, tipo, \emph{seq}, \emph{ack}, \emph{payload\_len},
autenticação de 32~bytes e \emph{checksum} CRC32), seguido de \emph{payload} de
até 4096~bytes. Os tipos de mensagem permanecem SYN, DATA, ACK e FIN. A
requisição HTTP é fragmentada em um ou mais pacotes SYN/DATA; a resposta HTTP
retorna pelo mesmo mecanismo. O \emph{timeout} R-UDP base é ajustado por cenário
(0,2~s, 0,5~s ou 1,0~s), coerente com o RTT emulado.

\subsection{Cenários de rede simulados}

A degradação foi aplicada com \texttt{tc qdisc} do pacote \texttt{iproute2},
executado no contêiner cliente pelo script \texttt{setup\_tc.sh}. Antes de cada
cenário, qualquer disciplina existente é removida; em seguida aplica-se o perfil
\texttt{netem} correspondente. A Tabela~\ref{tab:cenarios} resume os três
cenários avaliados.

\begin{table}[!htb]
  \centering
  \caption{Cenários de rede e comandos de simulação}
  \label{tab:cenarios}
  \begin{tabular}{cccc}
    \toprule
    Cenário & Perda & Atraso & Comando \\
    \midrule
    A & 0\%  & 10~ms  & \texttt{setup\_tc.sh A eth0} \\
    B & 5\%  & 50~ms  & \texttt{setup\_tc.sh B eth0} \\
    C & 10\% & 100~ms & \texttt{setup\_tc.sh C eth0} \\
    \bottomrule
  \end{tabular}
\end{table}

\subsection{Procedimento de testes}

O procedimento seguido em cada sessão obedeceu à sequência abaixo:

\begin{enumerate}
  \item Construção das imagens e inicialização dos contêineres com
        \texttt{docker compose up -d}, aguardando DNS, servidor e cliente.
  \item Geração dos arquivos em \texttt{www/files/} e limpeza de
        \texttt{logs/transfers.jsonl} para evitar mistura entre execuções.
  \item Para cada cenário A, B e C: aplicação do perfil \texttt{netem} no cliente;
        execução de dez repetições por combinação de modo (\texttt{tcp} ou
        \texttt{rudp}) e recurso (\texttt{/index.html},
        \texttt{/files/arquivo\_100kb.bin}, \texttt{/files/arquivo\_1mb.bin},
        \texttt{/files/arquivo\_10mb.bin}) via \texttt{run\_benchmark.sh}.
  \item Remoção das regras \texttt{tc} ao final dos três cenários.
  \item Geração dos gráficos estatísticos com \texttt{analysis/analyze.py}
        (Pandas e Matplotlib) e da taxa de erro com
        \texttt{analysis/gerar\_taxa\_erro.py} a partir do log do terminal.
  \item Capturas dedicadas com \texttt{tcpdump -i eth0} dentro do contêiner
        cliente, filtrando tráfego para \texttt{172.29.0.5} e
        \texttt{172.29.0.10}, durante execuções manuais representativas.
\end{enumerate}

No total, planejaram-se 240 transferências (3 cenários $\times$ 2 modos
$\times$ 4 arquivos $\times$ 10 repetições). Cada invocação de
\texttt{web\_client.py} registra separadamente \texttt{dns\_duration\_s},
\texttt{http\_duration\_s} e \texttt{total\_duration\_s}, permitindo isolar o
impacto da resolução de nomes do tempo de download HTTP.

\subsection{Validação cruzada}

A validação cruzada confronta métricas da aplicação Python com observações
diretas nas capturas Wireshark/\texttt{tshark}. A Tabela~\ref{tab:validacao}
resume a correspondência; os fluxos foram filtrados por
\texttt{ip.addr==172.29.0.20}, \texttt{udp.port==53}, \texttt{tcp.port==8080} e
\texttt{udp.port==8081}. As capturas dedicadas referem-se a execuções
representativas, distintas das 240 rodadas estatísticas do \emph{benchmark}.

\begin{table}[!htb]
  \centering
  \caption{Correspondência entre métricas da aplicação e medições de rede}
  \label{tab:validacao}
  \small
  \begin{tabularx}{\linewidth}{|p{3.2cm}|X|}
    \hline
    \textbf{Métrica (Python)} & \textbf{Medição no Wireshark/tshark} \\
    \hline
    \texttt{dns\_duration\_s} &
    Intervalo entre primeiro UDP:53 cliente$\rightarrow$DNS e resposta
    DNS$\rightarrow$cliente (\emph{frame.time\_relative}). \\
    \hline
    \texttt{http\_duration\_s} &
    TCP: SYN até último segmento da resposta HTTP.
    R-UDP: primeiro pacote UDP:8081 até FIN da sessão. \\
    \hline
    \texttt{bytes\_received} &
    Tamanho do corpo HTTP (\texttt{Content-Length}) ou soma dos
    \emph{payloads} DATA R-UDP após remoção dos cabeçalhos HTTP. \\
    \hline
    \texttt{retransmissions} &
    Contador interno do cliente R-UDP; no Wireshark, pacotes DATA repetidos
    com mesmo \emph{seq} ou SYN reenviado antes do ACK. \\
    \hline
    Ordem DNS $\rightarrow$ HTTP &
    Sequência UDP:53 completa antes do SYN TCP:8080 ou do primeiro UDP:8081. \\
    \hline
  \end{tabularx}
\end{table}

\section{Resultados Experimentais}

\subsection{Desempenho no cenário A}

Na ausência de perda, o TCP manteve 100\% de sucesso em todas as combinações.
Para o arquivo de 10~MB, a vazão média alcançou 126,0~Mbps (desvio de 9,7~Mbps),
enquanto o R-UDP estabilizou-se em torno de 3,1~Mbps. A diferença decorre do
Stop-and-Wait: cada \emph{chunk} de aproximadamente 4~kB aguarda ACK antes do
próximo envio, enquanto o TCP aproveita janela deslizante do \emph{kernel}.

O tempo médio de DNS ficou em 10~ms para ambos os modos, representando cerca de
14\% do tempo total no TCP para arquivos pequenos e menos de 2\% no arquivo de
10~MB. A Tabela~\ref{tab:vazao-a} resume os principais valores do cenário A.

\begin{table}[!htb]
  \centering
  \caption{Vazão média (Mbps) no cenário A}
  \label{tab:vazao-a}
  \begin{tabular}{lrr}
    \toprule
    Arquivo & TCP & R-UDP \\
    \midrule
    \texttt{index.html} (624~B) & 0,14 & 0,11 \\
    100~kB & 12,8 & 2,7 \\
    1~MB & 62,8 & 3,0 \\
    10~MB & 126,0 & 3,1 \\
    \bottomrule
  \end{tabular}
\end{table}

A Figura~\ref{fig:throughput-10mb} ilustra a diferença de vazão entre TCP e
R-UDP para o arquivo de 10~MB nos três cenários simulados.

\begin{figure}[!htb]
  \centering
  \includegraphics[width=0.92\textwidth]{throughput_arquivo_10mb_bin.png}
  \caption{Throughput médio para \texttt{arquivo\_10mb.bin} (10~MB)}
  \label{fig:throughput-10mb}
\end{figure}

O mesmo padrão aparece no arquivo de 1~MB (Figura~\ref{fig:throughput-1mb}),
com queda acentuada do R-UDP nos cenários B e C.

\begin{figure}[!htb]
  \centering
  \includegraphics[width=0.92\textwidth]{throughput_arquivo_1mb_bin.png}
  \caption{Throughput médio para \texttt{arquivo\_1mb.bin}}
  \label{fig:throughput-1mb}
\end{figure}

Para o arquivo de 100~kB, a Figura~\ref{fig:throughput-100kb} mostra vazão TCP
próxima de 13~Mbps no cenário A e degradação progressiva com a perda simulada.

\begin{figure}[!htb]
  \centering
  \includegraphics[width=0.92\textwidth]{throughput_arquivo_100kb_bin.png}
  \caption{Throughput médio para \texttt{arquivo\_100kb.bin}}
  \label{fig:throughput-100kb}
\end{figure}

A Figura~\ref{fig:throughput-index} apresenta o comportamento para
\texttt{index.html}, em que o tempo fixo de DNS e cabeçalhos HTTP domina o
cálculo de vazão aparente.

\begin{figure}[!htb]
  \centering
  \includegraphics[width=0.92\textwidth]{throughput_index_html.png}
  \caption{Throughput médio para \texttt{index.html} (624~bytes)}
  \label{fig:throughput-index}
\end{figure}

\subsection{Impacto da perda nos cenários B e C}

O TCP manteve taxa de erro nula nos cenários B e C, embora a vazão do arquivo de
10~MB caísse de 126~Mbps (A) para 41~Mbps (B) e 24~Mbps (C). Já o R-UDP
apresentou falhas expressivas a partir do cenário B para arquivos maiores que
\texttt{index.html}: 90\% de erro no arquivo de 100~kB, 100\% em 1~MB e 10~MB.
No cenário C, mesmo \texttt{index.html} registrou 30\% de falha (7 sucessos em
10 tentativas). A Tabela~\ref{tab:erro} consolida a taxa de erro inferida do
log de execução do \emph{benchmark}.

\begin{table}[!htb]
  \centering
  \caption{Taxa de erro (\%) por cenário, modo e arquivo}
  \label{tab:erro}
  \small
  \begin{tabular}{clcc}
    \toprule
    Cen. & Arquivo & TCP & R-UDP \\
    \midrule
    \multirow{4}{*}{B}
      & \texttt{index.html} & 0 & 0 \\
      & 100~kB & 0 & 90 \\
      & 1~MB & 0 & 100 \\
      & 10~MB & 0 & 100 \\
    \midrule
    \multirow{4}{*}{C}
      & \texttt{index.html} & 0 & 30 \\
      & 100~kB & 0 & 100 \\
      & 1~MB & 0 & 100 \\
      & 10~MB & 0 & 100 \\
    \bottomrule
  \end{tabular}
\end{table}

A Figura~\ref{fig:taxa-erro} apresenta visualmente a taxa de falha para todas
as combinações testadas, incluindo barras de 100\% para R-UDP em arquivos
grandes nos cenários degradados.

\begin{figure}[!htb]
  \centering
  \includegraphics[width=0.92\textwidth]{taxa_erro.png}
  \caption{Taxa de erro por cenário, modo de transporte e tamanho de arquivo}
  \label{fig:taxa-erro}
\end{figure}

Consultas DNS com retransmissão apareceram nos logs com duração próxima de
1,05~s (cenário B) e 2,10~s (cenário C). A Figura~\ref{fig:dns-http} contrasta
o tempo médio de DNS e de HTTP por cenário e modo de transporte.

\begin{figure}[!htb]
  \centering
  \includegraphics[width=0.92\textwidth]{dns_vs_http_tempo.png}
  \caption{Tempo médio de DNS e de HTTP por cenário e modo}
  \label{fig:dns-http}
\end{figure}

Esses picos concentraram-se em menos de 10\% das execuções TCP, mas tornaram o
DNS proporcionalmente mais relevante em páginas pequenas. A
Figura~\ref{fig:duracao-box} exibe a dispersão do tempo total de carregamento
(DNS + HTTP) em cada combinação.

\begin{figure}[!htb]
  \centering
  \includegraphics[width=0.92\textwidth]{duracao_total_boxplot.png}
  \caption{Distribuição do tempo total de carregamento (DNS + HTTP)}
  \label{fig:duracao-box}
\end{figure}

\subsection{Validação Wireshark}

A Figura~\ref{fig:fluxo-geral} exibe a listagem filtrada por
\texttt{ip.addr==172.29.0.20}, evidenciando consulta UDP ao DNS seguida de
\emph{handshake} TCP e requisição GET.

\begin{figure}[!htb]
  \centering
  \includegraphics[width=0.92\textwidth]{"ip.addr==172.29.0.20.png"}
  \caption{Fluxo completo DNS $\rightarrow$ HTTP/TCP filtrado pelo cliente}
  \label{fig:fluxo-geral}
\end{figure}

A Figura~\ref{fig:transicao-dns-tcp} amplia os \emph{frames} 1 a 6: a resposta
DNS encerra em 0,001~s e o SYN TCP parte em 0,013~s, intervalo de
aproximadamente 12~ms.

\begin{figure}[!htb]
  \centering
  \includegraphics[width=0.92\textwidth]{"frame.number >= 1 && frame.number <= 6.png"}
  \caption{Transição entre resposta DNS e início do \emph{handshake} TCP}
  \label{fig:transicao-dns-tcp}
\end{figure}

A Figura~\ref{fig:get} detalha os cabeçalhos da requisição GET, incluindo
\texttt{Host} e \texttt{X-Custom-Auth}.

\begin{figure}[!htb]
  \centering
  \includegraphics[width=0.88\textwidth]{1.1.png}
  \caption{Detalhamento da requisição HTTP GET}
  \label{fig:get}
\end{figure}

A Figura~\ref{fig:200ok} apresenta a resposta \texttt{HTTP/1.1 200 OK} com
\texttt{Content-Type}, \texttt{Content-Length} e \texttt{X-Custom-Auth}.

\begin{figure}[!htb]
  \centering
  \includegraphics[width=0.88\textwidth]{"200 OK.png"}
  \caption{Resposta HTTP/1.1 200 OK com cabeçalhos exigidos}
  \label{fig:200ok}
\end{figure}

A Figura~\ref{fig:dns-rudp} mostra a sequência DNS seguida de tráfego UDP na
porta 8081, sem \emph{handshake} TCP.

\begin{figure}[!htb]
  \centering
  \includegraphics[width=0.92\textwidth]{"udp.port == 53 || udp.port == 8081.png"}
  \caption{Consulta DNS e início da sessão R-UDP (porta 8081)}
  \label{fig:dns-rudp}
\end{figure}

O diagrama de fluxo da Figura~\ref{fig:flow} resume visualmente a ordem das
interações entre cliente, DNS e servidor web.

\begin{figure}[!htb]
  \centering
  \includegraphics[width=0.92\textwidth]{flow.png}
  \caption{Diagrama de fluxo gerado pelo Wireshark (\emph{Flow Graph})}
  \label{fig:flow}
\end{figure}

No cenário B com R-UDP, a Figura~\ref{fig:rudp-b} registra múltiplos pacotes
UDP:8081 em sequência Stop-and-Wait, compatível com retransmissões observadas
durante o \emph{benchmark}. O Wireshark rotula as consultas DNS simplificadas
como \emph{Malformed Packet}, pois o formato adotado difere do RFC~1035; a
ordem temporal, contudo, confirma o comportamento esperado do sistema.

\begin{figure}[!htb]
  \centering
  \includegraphics[width=0.92\textwidth]{"cenario_B_rudp.pcap-Filtro: udp.port==8081.png"}
  \caption{Tráfego R-UDP no cenário B com 5\% de perda}
  \label{fig:rudp-b}
\end{figure}

\section{Discussão das Perguntas Obrigatórias}

\subsection{Perda simulada, DNS e download HTTP}

A perda introduzida pelo \texttt{netem} afetou o DNS de forma pontual, mas
mensurável. Como o UDP não retransmite automaticamente, o cliente DNS recorreu
ao temporizador de aplicação: consultas bem-sucedidas mantiveram cerca de 50~ms
(B) e 100~ms (C), alinhadas ao atraso configurado, enquanto falhas isoladas
elevaram o tempo para 1~s ou 2~s conforme o número de \emph{retries}. O impacto
relativo foi maior em objetos pequenos, nos quais o DNS pode representar até 18\%
do tempo total TCP, conforme a Figura~\ref{fig:dns-http}.

No download HTTP via TCP, a pilha do \emph{kernel} absorveu perdas e
retransmissões de forma transparente, preservando 100\% de sucesso. Já o HTTP
sobre R-UDP degradou drasticamente: a partir de 5\% de perda, transferências
acima de 100~kB praticamente não completaram dentro do limite de tentativas, o
que se reflete na Figura~\ref{fig:taxa-erro}. Confirmou-se, portanto, que o
DNS exige \emph{timeout} explícito na aplicação, enquanto HTTP/TCP delega
confiabilidade ao transporte nativo.

\subsection{Overhead dos cabeçalhos HTTP}

Comparando com o protocolo textual customizado do trabalho anterior, a camada
HTTP adiciona linha de requisição, múltiplos cabeçalhos e separador
\texttt{\textbackslash r\textbackslash n\textbackslash r\textbackslash n}. Para
\texttt{index.html} (624~bytes), o tempo total TCP no cenário A foi 36~ms, dos
quais cerca de 31~ms correspondem a DNS e negociação HTTP, ou seja, 86\% do
tempo gasto em \emph{overhead} fixo. Para o arquivo de 10~MB (669~ms total),
esse \emph{overhead} cai para aproximadamente 5\%.

A Figura~\ref{fig:throughput-index} mostra vazões abaixo de 0,2~Mbps para a
página HTML, enquanto a Figura~\ref{fig:throughput-10mb} evidencia dezenas a
centenas de Mbps no TCP para o arquivo grande. Essa diferença visual confirma
que o custo dos cabeçalhos HTTP concentra-se nas transferências de poucos bytes
e torna-se marginal quando o payload cresce.

\subsection{Ordem do fluxo na arquitetura em camadas}

As capturas confirmam a sequência prevista pela arquitetura da Internet: camada
de aplicação DNS (UDP/53) precedendo a sessão HTTP. Nos \emph{frames} analisados,
observou-se consulta DNS (\texttt{172.29.0.20} $\rightarrow$ \texttt{172.29.0.5}),
resposta imediata, SYN TCP para \texttt{172.29.0.10:8080}, \emph{handshake}
completo, GET \texttt{/index.html}, resposta \texttt{200 OK} e encerramento FIN,
conforme Figuras~\ref{fig:fluxo-geral} e~\ref{fig:transicao-dns-tcp}.

No modo R-UDP, após a mesma resolução DNS, o primeiro pacote de dados aparece
em UDP/8081, sem \emph{handshake} TCP, como ilustram as
Figuras~\ref{fig:dns-rudp},~\ref{fig:flow} e~\ref{fig:rudp-b}.

\section{Conclusão}

Este trabalho integrou resolução DNS local, miniservidor HTTP/1.1 e transporte
alternável entre TCP e R-UDP em ambiente Docker instrumentado. Os experimentos
repetidos em três cenários de rede quantificaram diferenças substanciais: o TCP
nativo manteve confiabilidade total e vazões superiores, enquanto o R-UDP
Stop-and-Wait mostrou-se adequado apenas no cenário A sem perda, falhando de
forma massiva quando a taxa de perda atingiu 5\% ou 10\% em arquivos maiores.

As capturas Wireshark corroboraram a ordem DNS $\rightarrow$ HTTP e permitiram
inspecionar cabeçalhos exigidos pelo enunciado. Como continuidade, sugere-se
incorporar registro sistemático de falhas no arquivo JSONL, ampliar \emph{timeouts}
R-UDP em cenários degradados ou adotar janela deslizante para reduzir sensibilidade
à perda. O repositório contém código, \emph{scripts} de \emph{benchmark}, gráficos
e arquivos \texttt{.pcap} utilizados nesta análise.

\bibliographystyle{sbc}
\bibliography{sbc-template}

\end{document}
